\begingroup
\renewcommand{\arraystretch}{1.45}
\fancyhead[L]{\small ARITHMETIC–GEOMETRIC MEAN.}
\Needspace{28\baselineskip}
\AGMArticle{21}

[If the equations between the derivatives of $\alpha,\allowbreak \beta,\allowbreak \gamma,\allowbreak \delta$ with respect to the quantity $y$ as independent variable are formed, the same coefficients of these derivatives must arise as in the equations for the differentials with respect to $\eta$. Since the latter equations can be put into a form that determines the ratios between the derivatives of quotients, the coefficients mentioned must also be inversely proportional to these derivatives. Replacing them by the reciprocal values shows that the remaining terms of those equations can also be represented by such derivatives, in such a way that the expression of the form
\Needspace{10\baselineskip}
\[
\dfrac{\left[\dfrac{\mathrm{d}\log\dfrac{\delta}{\beta}}{\mathrm{d}\log y}\right]}
{\left[\dfrac{\mathrm{d}\log\dfrac{\delta}{\beta}}{\mathrm{d}\log\eta}\right]}
-\tfrac12\dfrac{\mathrm{d}\log\dfrac{\alpha\beta}{xx}}{\mathrm{d}\log\eta}
\]
does not change its value for any common index of the symbols $\alpha,\allowbreak \beta,\allowbreak \gamma,\allowbreak \delta,\allowbreak x,\allowbreak y,\allowbreak \eta$, in whatever order the quantities $\alpha,\allowbreak \beta,\allowbreak \gamma,\allowbreak \delta$ are interchanged. Comparing this expression, which refers to $\frac{\delta}{\beta},\allowbreak \frac{\alpha\gamma}{xx},\allowbreak \eta,\allowbreak y$, with that for $\frac{\delta'}{\beta'}=\frac{\gamma\delta}{\alpha\beta},\allowbreak \frac{\alpha'\gamma'}{x'x'},\allowbreak \eta'=\eta\eta,\allowbreak y'=yy$ and observing that, by Article 19,
\Needspace{14\baselineskip}
\[
\begin{gathered}\left[\dfrac{\mathrm{d}\log\dfrac{\alpha}{\beta}}{\mathrm{d}\log\eta}\right]^2
=\dfrac{\mathrm{d}\log\dfrac{\gamma'}{\beta'}}{\mathrm{d}\log\eta}\cdot
\dfrac{\mathrm{d}\log\dfrac{\alpha'}{\beta'}}{\mathrm{d}\log\eta},\\[9pt]
\left[\dfrac{\mathrm{d}\log\dfrac{\delta}{\gamma}}{\mathrm{d}\log\eta}\right]^2
=\dfrac{\mathrm{d}\log\dfrac{\delta'}{\alpha'}}{\mathrm{d}\log\eta}\cdot
\dfrac{\mathrm{d}\log\dfrac{\delta'}{\gamma'}}{\mathrm{d}\log\eta}\end{gathered}
\]
shows that the value of that expression must also be independent of the index of the symbols $\alpha,\allowbreak \beta,\allowbreak \gamma,\allowbreak \delta,\allowbreak x,\allowbreak y,\allowbreak \eta$ and, as follows from the limits, must equal zero.
\Needspace{22\baselineskip}
\SourceLeaf{402}

To determine the derivative with respect to $y$ of each of the four quantities $\alpha,\allowbreak \beta,\allowbreak \gamma,\allowbreak \delta$, it will first be necessary to transform that expression so that it contains only two of the quantities. This can be done, for example, with the aid of the second derivatives of the quotients with respect to $\eta$ determined in the preceding article, so that all expressions
\Needspace{12\baselineskip}
\[
\begin{aligned}&\dfrac{\mathrm d\log\dfrac{\alpha}{\beta}}{\mathrm d\log y}
-\dfrac{\mathrm d\mathrm d\log\dfrac{\alpha}{\beta}}{(\mathrm d\log\eta)^2}
\\[8pt]&{}-\dfrac12\dfrac{\mathrm d\log\dfrac{\alpha\beta}{xx}}{\mathrm d\log\eta}
\cdot\dfrac{\mathrm d\log\dfrac{\alpha}{\beta}}{\mathrm d\log\eta}\end{aligned}
\]
vanish, in whatever order $\alpha,\allowbreak \beta,\allowbreak \gamma,\allowbreak \delta$ are interchanged and whatever common index is given to them and to the symbols $y,\allowbreak x,\allowbreak \eta$. If $\frac{\alpha}{\beta}$ is replaced by $\frac{b'\beta'}{\beta\beta}$, $\alpha\beta$ by $b'\beta'$ and, according to the preceding article,
\Needspace{14\baselineskip}
\[
\begin{gathered}\dfrac{\mathrm d\log\dfrac{b'}{k}}{2\mathrm d\log y}
\quad\text{by}\\[9pt]
\dfrac{\mathrm d\mathrm d\log\dfrac{\beta'}{x}}{(2\mathrm d\log\eta)^2}
+\dfrac12\left[\dfrac{\mathrm d\log\dfrac{\beta\beta}{\beta'}}{\mathrm d\log\eta}\right]^2\end{gathered}
\]
it is seen that the expression
\Needspace{12\baselineskip}
\[
\begin{aligned}&\dfrac{\mathrm d\log\dfrac{\beta}{x}}{\mathrm d\log y}
-\dfrac{\mathrm d\mathrm d\log\dfrac{\beta}{x}}{(\mathrm d\log\eta)^2}
\\[8pt]&{}-\dfrac12\left[\dfrac{\mathrm d\log\dfrac{\beta}{x}}{\mathrm d\log\eta}\right]^2\end{aligned}
\]
remains unchanged in value, not only, as just found, when $\beta$ is interchanged with $\alpha$ or $\gamma$ or $\delta$, but also when these quantities, together with $y,\allowbreak x,\allowbreak \eta$, are given any common index. For an index that increases steadily, the value of the expression vanishes, and multiplication by $\frac12\surd {\frac{\beta}{x}}$ therefore gives
\Needspace{10\baselineskip}
\[
\dfrac{\mathrm d\surd {\dfrac{\alpha}{x}}}{\mathrm d\log y}
-\dfrac{\mathrm d\mathrm d\surd {\dfrac{\alpha}{x}}}{(\mathrm d\log\eta)^2}=0
\]
as a partial differential equation for each of the quantities $\alpha,\allowbreak \beta,\allowbreak \gamma,\allowbreak \delta$.]

\Needspace{26\baselineskip}
\AGMArticle{22}
[The values of all terms of the algorithm depend on four quantities, which may be taken as $a,\allowbreak b,\allowbreak \alpha,\allowbreak \beta$, but also as $k,\allowbreak y,\allowbreak x,\allowbreak \eta$, since these can vary independently of one another. All equations between the terms of the progressions can be put into a form that is homogeneous both in $a,\allowbreak b,\allowbreak c,\allowbreak .\;. a^{\mathrm n},\allowbreak b^{\mathrm n},\allowbreak c^{\mathrm n},\allowbreak k$ and in


\SourceLeaf{403}

\Needspace{10\baselineskip}
\[
\begin{gathered}\dfrac{\alpha}{k},\ \dfrac{\beta}{k},\ \dfrac{\gamma}{k},\ \dfrac{\delta}{k}\ .\;.\;.\\[8pt]\
{}^{2^n}\!\surd {\dfrac{\alpha^{\mathrm n}}{k}},\
{}^{2^n}\!\surd {\dfrac{\beta^{\mathrm n}}{k}},\
{}^{2^n}\!\surd {\dfrac{\gamma^{\mathrm n}}{k}},\
{}^{2^n}\!\surd {\dfrac{\delta^{\mathrm n}}{k}},\ \dfrac{x}{k}\end{gathered}
\]
The ratios between the latter therefore depend neither on $k$ nor on $x$, but only on $y$ and $\eta$; we can therefore set
\Needspace{14\baselineskip}
\[
\begin{gathered}\alpha=x P(y,\eta)^2,\\[5pt]
\beta=x Q(y,\eta)^2,\\[5pt]
\gamma=x R(y,\eta)^2,\\[5pt]
\delta=x S(y,\eta)^2\end{gathered}
\]
and it then follows from the considerations in Article 18 that these functions $P,\allowbreak \ Q,\allowbreak \ R,\allowbreak \ S$ remain the same when $\alpha,\allowbreak \beta,\allowbreak \gamma,\allowbreak \delta,\allowbreak x,\allowbreak y,\allowbreak \eta$ are given any index. \Needspace{16\baselineskip}
From the same article, using also the notation employed in Article 15, it follows for $a,\allowbreak b. c,\allowbreak y$ bearing the same index that
\[
\begin{aligned}
P(y,1)&=py=\surd {\dfrac ak}=\surd {\dfrac a{\mathrm M(a,b)}}\\
Q(y,1)&=qy=\surd {\dfrac bk}\\
R(y,1)&=ry=\surd {\dfrac ck}\\
S(y,1)&=0
\end{aligned}
\]
and, further, as equations corresponding to those by which the quantities $\gamma$ and $\delta$ were determined:
\Needspace{24\baselineskip}
\[
\begin{gathered}
qqrr(qqRR-rrQQ)\\[4pt]=rrpp(ppRR-rrPP)\\
=ppqq(qqPP-ppQQ)\\[4pt]=ppqqrrSS\\[4pt]
\dfrac{PP+QQ}{p(yy)}=\dfrac{RR-SS}{r(yy)}.\\[9pt]
\dfrac{PP-QQ}{r(yy)}=\dfrac{RR+SS}{p(yy)}\\[4pt]
\dfrac{PP+RR}{p(\surd  y)}=\dfrac{QQ+SS}{q(\surd  y)},\\[9pt]
\dfrac{PP-RR}{q(\surd  y)}=\dfrac{QQ-SS}{p(\surd  y)}
\end{gathered}
\]
\Needspace{12\baselineskip}
and as equations representing the law by which the algorithm is formed:
\[
\begin{aligned}
2p(yy)\cdot P(yy,\eta\eta)&=PP+QQ,\\[5pt]
2p(yy)\cdot R(yy,\eta\eta)&=RR+SS\\[5pt]
2r(yy)\cdot R(yy,\eta\eta)&=PP-QQ,\\[5pt]
2r(yy)\cdot P(yy,\eta\eta)&=RR-SS\\[5pt]
q(yy)\cdot Q(yy,\eta\eta)&=PQ,\\[5pt]
q(yy)\cdot S(yy,\eta\eta)&=RS
\end{aligned}
\]
where all functions beside which no argument is written refer to the simple $y$ and $\eta$.

\Needspace{22\baselineskip}
\SourceLeaf{404}

To calculate the values of the functions that belong to given values of $\frac ab=\frac{pp}{qq}$ and $U$ or $V$, the following equations are obtained from \textsc{Gauss}'s formula for $H$ or $x$, if, for brevity,
\Needspace{10\baselineskip}
\[
\begin{aligned}&(a^{\mathrm n}\sin 2^{\mathrm n} U^{\mathrm n})^2+(b^{\mathrm n}\cos 2^{\mathrm n} U^{\mathrm n})^2\\[5pt]&=b^{\mathrm n} b^{\mathrm n}\Delta^{\mathrm n}\Delta^{\mathrm n}\end{aligned}
\]
is set:
\Needspace{18\baselineskip}
\[
\begin{aligned}
Q&=q\cdot\surd \Delta\cdot{}^{4}\!\surd {\Delta'}\cdot{}^{8}\!\surd {\Delta''}\cdot{}^{16}\!\surd {\Delta'''}.\;.\;.\\[9pt]
P&=\dfrac pq Q\dfrac1\Delta\\[9pt]
R&=\dfrac rq Q\dfrac{\cos U}{\Delta}\\[9pt]
S&=\dfrac pr Q\dfrac{\sin U}{\Delta}i
\end{aligned}
\]

\Needspace{20\baselineskip}
The determination of the quantities $k,\allowbreak y,\allowbreak x,\allowbreak \eta$ as limits shows that, with a suitable choice of the signs of the functions $P,\allowbreak Q,\allowbreak R,\allowbreak S$, which is still permissible here, for values of $y,\allowbreak y\eta,\allowbreak \frac1\eta$ decreasing to zero, the expressions
\Needspace{16\baselineskip}
\[
\begin{gathered}P(y,\eta),\\[6pt] Q(y,\eta).\\[6pt]
\dfrac{R(y,\eta)}{y^{\frac14}\eta^{\frac12}},\\[6pt]
\dfrac{S(y,\eta)}{y^{\frac14}\eta^{\frac12}}\end{gathered}
\]
\Needspace{16\baselineskip}
approach the limit unity; but for a steadily decreasing $y$ and a finite real $u$ the expressions
\[
\begin{gathered}P(y,e^{2iu}),\\[6pt] Q(y,e^{2iu}),\\[6pt]
\dfrac{R(y,e^{2iu})}{2y^{\frac14}\cos u},\\[6pt]
\dfrac{S(y,e^{2iu})}{i2y^{\frac14}\sin u}\end{gathered}
\]
have that limit.

\Needspace{16\baselineskip}
The functions $P,\allowbreak Q,\allowbreak R,\allowbreak iS$ have real values for a complex $\eta^{\frac12}$ of the form $e^{iu}$ and remain unchanged except for $iS$, which merely changes sign, when $u$ is changed into $-u$; thus
\Needspace{10\baselineskip}
\[
\begin{aligned}\dfrac{P(y,\dfrac1\eta)}P
&=\dfrac{Q(y,\dfrac1\eta)}Q
\\[7pt]&=\dfrac{R(y,\dfrac1\eta)}R
\\[7pt]&=\dfrac{-S(y,\dfrac1\eta)}S=1
\text{]}\end{aligned}
\]

\clearpage
\endgroup
